\chapter{Credit Indices and Tranches}\label{ch:m2:credit-indices-and-tranches}

In the first quarter of 2012 the Chief Investment Office of JPMorgan Chase tripled
the net notional of a portfolio of credit indices and \omterm{def:m2:credit-indices-and-tranches:tranche}{tranches}, created in 2007 as a
hedge against credit losses, to about USD~157 billion. A large part was protection
sold on the ten-year contract of one series of the North American investment-grade
index, issued in 2007; at one point the net position in that series was USD~82
billion, 10 to 15 days of the whole market's trading volume.
Its traders came to believe that the market knew their positions and was pricing
against them. The bank disclosed a loss of about USD~2 billion in May; by the end of
the year the losses stood at USD~6.2 billion, and it later admitted the facts in a
settlement with the SEC. This chapter explains the instruments: credit indices and how
they roll, the gap between an index and its constituents, the \omterm{def:m2:credit-indices-and-tranches:tranche}{tranches} that slice an
index's losses, and options on indices.

\section{The index families and their rolls}

\begin{definition}[Credit index, index series]\label{def:m2:credit-indices-and-tranches:index}
A \emph{credit index}\index{credit index} is a standardised \omterm{def:m2:credit-default-swaps:cds}{credit default swap} on an
equally weighted basket of reference entities, with one coupon and one maturity; a
credit event of a constituent is settled on that constituent's share of the notional,
and the index continues on the others. An \emph{index series}\index{index series} is
one fixed list of constituents, published on a roll date; a new series, with an updated
list, is published every six months, and the most recent is the on-the-run series.
\end{definition}

Two families dominate: CDX, North American and emerging-market indices, and iTraxx,
European and Asian ones, each with investment-grade and high-yield members. The
North American investment-grade index holds 125 of the most liquid investment-grade
names, with equal weights, and rolls on 20 March and 20 September; a name that has
been downgraded, stopped trading or been taken over is replaced in the new series,
while the old series keeps its list and trades on, less and less, off the run. The
index pays a fixed coupon of 100 basis points quarterly and is quoted as a spread, with
the upfront settled as for a single name (chapter 23). The 2007 series the bank's
office traded in 2012 was the ninth; the series issued in September 2012 was the
nineteenth.

An index lets an investor trade credit as a whole: one trade in one liquid
instrument buys or sells protection on a whole market, far more cheaply than 125
single names. It can hedge a book of single names or a bond portfolio, or express a
view on the economy.

\begin{dated}{2026-09}{Index contract facts}\label{dat:m2:credit-indices-and-tranches:indices}
CDX North American investment grade: 125 names, equal weights, rolls on 20 March and
20 September, tenors 1 to 10 years, 100 basis point coupon, cleared at ICE Clear
Credit. Index options: European-style, physically settled, on the five-year CDX
investment-grade and high-yield, iTraxx Europe and iTraxx Crossover indices, up to
nine months to expiry, cleared at ICE.
\end{dated}

\section{Index versus constituents}

An index is a basket of CDS, so it has a value implied by the prices of those CDS.
The index's own quote, set by trading in the index, need not agree.

\begin{definition}[Intrinsic spread, index skew]\label{def:m2:credit-indices-and-tranches:intrinsic}
The \emph{intrinsic spread}\index{intrinsic spread} of an index is the spread that,
converted with the index's coupon and conventions, gives the weighted average of its
constituents' upfronts at that coupon. The \emph{index skew}\index{index skew} is the
index's quoted spread minus its intrinsic spread.
\end{definition}

\begin{proposition}[Intrinsic spread as a weighted average]\label{prop:m2:credit-indices-and-tranches:intrinsic}
If constituent $i$ has weight $w_i$, spread $s_i$ and risky annuity $A_i$, the index
with coupon $c$ has intrinsic upfront $\sum_i w_i (s_i - c) A_i$, and, to first
order, \omterm{def:m2:credit-indices-and-tranches:intrinsic}{intrinsic spread}
\[
s^{\text{int}} \approx \frac{\sum_i w_i A_i s_i}{\sum_i w_i A_i},
\]
which is below the simple average of the spreads, since names with higher spreads have
smaller annuities.
\end{proposition}

\begin{proof}
The upfront of each constituent is $(s_i - c)A_i$ by chapter 23, and the index's
protection and premium legs are the weighted sums of the constituents', so its
upfront is their weighted sum. Writing it as $(s^{\text{int}} - c)A$ with $A \approx
\sum_i w_i A_i$ gives the formula. A higher spread means a higher hazard rate and so
a shorter expected life, hence a smaller annuity and a smaller weight.
\end{proof}

\begin{example}[A stylised index]\label{ex:m2:credit-indices-and-tranches:intrinsic}
Take 125 names whose five-year spreads are spread log-normally around a median of 55
basis points, from 6.6 to 459, with a simple average of 75.3. With a flat rate of 4\%
and a recovery of 40\%, their average upfront at the 100 basis point coupon is
$-1.157\%$, which converts to an \omterm{def:m2:credit-indices-and-tranches:intrinsic}{intrinsic spread} of 73.6 basis points; the
annuity-weighted average of the spreads is also 73.6. If the index trades at 65.6, its
skew is $-8$ basis points (\cref{fig:m2:credit-indices-and-tranches:constituents}).
\end{example}

\begin{omfigure}
\begin{tikzpicture}
\begin{axis}[width=11cm,height=5cm,
  xlabel={constituent, sorted by spread},ylabel={five-year spread (bp)},
  tick label style={font=\small},label style={font=\small},
  xmin=0,xmax=126,ymin=0,ymax=300,grid=major,grid style={gray!25},
  legend columns=2,legend style={font=\footnotesize,at={(0.5,-0.3)},anchor=north,draw=none,fill=none,/tikz/every even column/.append style={column sep=8pt}},
  table/col sep=comma]
\addplot[omDef,only marks,mark=*,mark size=1pt] table[x=rank,y=spread]{figdata/markets-2/24-credit-indices-and-tranches/constituents.csv};
\addplot[gray,thick,dashed] coordinates {(0,75.3) (126,75.3)};
\addplot[omThm,very thick] coordinates {(0,73.6) (126,73.6)};
\addplot[omProp,very thick,dotted] coordinates {(0,65.6) (126,65.6)};
\legend{constituents,simple average 75.3,intrinsic 73.6,index quote 65.6}
\end{axis}
\end{tikzpicture}

\omcaption{The constituents of a stylised 125-name investment-grade index, sorted by
five-year spread (the two widest, at 335 and 459 basis points, are off the chart). The
\omterm{def:m2:credit-indices-and-tranches:intrinsic}{intrinsic spread} is below the simple average because wide names carry smaller
annuities; an index quoted at 65.6 trades 8 basis points through its intrinsic value.
Illustrative; data: the chapter's tutorial.}\label{fig:m2:credit-indices-and-tranches:constituents}
\end{omfigure}

The skew is an arbitrage in principle: buy protection on the index and sell it on the
125 names, and the position pays nothing on a default, since each default is settled
on both sides, and gains the skew when it closes. In practice the skew persists and
can widen. Trading 125 single names costs far more than one index trade; single names
are much less liquid than the index; the trade ties up margin and capital for as long
as the skew stays open; and a large, persistent buyer or seller of the index can push
it further away before it comes back. A skew is where arbitrage capital and one
side's demand meet. A very large seller of index protection pushes the index below its
constituents, and the investors who trade the skew are its natural counterparties; in
2012 the bank's traders came to believe that those counterparties knew their positions.

\begin{omfigure}
\begin{tikzpicture}
\begin{axis}[width=11cm,height=4.8cm,
  xlabel={index skew at exit (bp)},ylabel={P\&L (USD millions)},
  tick label style={font=\small},label style={font=\small},
  xmin=-24,xmax=8,ymin=-8,ymax=6,grid=major,grid style={gray!25},
  legend columns=2,legend style={font=\footnotesize,at={(0.5,-0.3)},anchor=north,draw=none,fill=none,/tikz/every even column/.append style={column sep=8pt}},
  table/col sep=comma]
\addplot[omDef,very thick] table[x=skew,y=pnl]{figdata/markets-2/24-credit-indices-and-tranches/skewpnl.csv};
\addplot[omProp,thick,dashed] coordinates {(-24,1.53) (8,1.53)};
\addplot[only marks,mark=*,mark size=1.8pt,black,forget plot] coordinates {(-8,0) (0,3.55) (-20,-5.37)};
\legend{gross P\&L,round-trip trading cost}
\end{axis}
\end{tikzpicture}

\omcaption{P\&L of a USD~1 billion skew trade entered at a skew of $-8$ basis points
(buy index protection, sell protection on the constituents) against the skew at exit,
constituents unchanged. The trade gains USD~3.55 million if the skew closes and loses
USD~5.37 million if it widens to $-20$ first; the dashed line is the cost of entering
and leaving both legs. Illustrative; data: the chapter's weekend
problem.}\label{fig:m2:credit-indices-and-tranches:skewpnl}
\end{omfigure}

\section{Tranches}

\begin{definition}[Tranche, attachment and detachment points]\label{def:m2:credit-indices-and-tranches:tranche}
A \emph{tranche}\index{tranche} of an index is a contract that covers only the
portfolio losses between two levels, expressed as fractions of the index notional:
losses start to reduce it at the \emph{attachment point}\index{attachment point} $a$
and have wiped it out at the \emph{detachment point}\index{detachment point} $d$. For a
portfolio loss $L$, the tranche loses $\min\bigl(\max(L - a, 0),\, d - a\bigr)$, a fraction
$\min(\max(L-a,0), d-a)/(d-a)$ of its own notional.
\end{definition}

The \omterm{def:m2:credit-indices-and-tranches:tranche}{tranches} of an index share its losses in order
(\cref{fig:m2:credit-indices-and-tranches:stack}): on the North American
investment-grade index the equity \omterm{def:m2:credit-indices-and-tranches:tranche}{tranche}, from 0 to 3\%, takes the first losses and
earns the highest premium; the mezzanine, from 3 to 7\%, is hit only after the equity
is gone; senior and super-senior \omterm{def:m2:credit-indices-and-tranches:tranche}{tranches} above them lose only in a catastrophe (in
this chapter's example, from 7 to 15\% and from 15 to 100\%). With a recovery of 40\%, each default of a 125-name index
costs 0.48\% of the notional, so the equity \omterm{def:m2:credit-indices-and-tranches:tranche}{tranche} is wiped out by seven defaults.
The expected losses of all the \omterm{def:m2:credit-indices-and-tranches:tranche}{tranches}, weighted by their widths, add up to the
expected loss of the whole portfolio, since together they cover it once. What a
\omterm{def:m2:credit-indices-and-tranches:tranche}{tranche} does not share with the index is its dependence on how defaults cluster.

\begin{omfigure}
\centering
\begin{tikzpicture}[font=\small]
\fill[omDef!35] (0,0) rectangle (3.8,0.6);
\fill[omThm!30] (0,0.6) rectangle (3.8,1.4);
\fill[omProp!25] (0,1.4) rectangle (3.8,2.6);
\fill[gray!15] (0,2.6) rectangle (3.8,5.2);
\draw[thick] (0,0) rectangle (3.8,5.2);
\foreach \y in {0.6,1.4,2.6} \draw[thick] (0,\y) -- (3.8,\y);
\node at (1.9,0.3) {equity 0--3\%};
\node at (1.9,1.0) {mezzanine 3--7\%};
\node at (1.9,2.0) {senior 7--15\%};
\node at (1.9,3.9) {super senior 15--100\%};
\node[anchor=west,font=\footnotesize] at (4.0,0) {0\%};
\node[anchor=west,font=\footnotesize] at (4.0,0.6) {3\%};
\node[anchor=west,font=\footnotesize] at (4.0,1.4) {7\%};
\node[anchor=west,font=\footnotesize] at (4.0,2.6) {15\%};
\node[anchor=west,font=\footnotesize] at (4.0,5.2) {100\%};
\draw[-{Stealth[length=2.5mm]},very thick,omDef] (-1.2,0) -- node[left,align=right,font=\footnotesize] {portfolio\\losses fill\\from the\\bottom} (-1.2,2.4);
\node[align=left,font=\footnotesize,anchor=west] at (5.4,0.3) {first loss, highest premium;\\wiped out by 7 defaults};
\node[align=left,font=\footnotesize,anchor=west] at (5.4,3.9) {loses only if more than\\31 defaults occur};
\end{tikzpicture}

\omcaption{The \omterm{def:m2:credit-indices-and-tranches:tranche}{tranches} of this chapter's example (heights not to scale). Portfolio losses
fill the stack from the bottom: the equity \omterm{def:m2:credit-indices-and-tranches:tranche}{tranche} takes the first 3\% of losses, the
mezzanine the next 4\%, and so on; with a 40\% recovery each of 125 defaults costs
0.48\% of the notional. Schematic.}\label{fig:m2:credit-indices-and-tranches:stack}
\end{omfigure}

The standard model of that dependence is the one-factor Gaussian model. Each name
defaults before the horizon if $\sqrt\rho\,Z + \sqrt{1-\rho}\,\varepsilon_i <
N^{-1}(p)$, with $Z$ a common factor, $\varepsilon_i$ independent standard normals, $p$
the probability of default to the horizon and $\rho$ the correlation between names.
Given $Z$, defaults are independent, and in a large pool the fraction that default is
their conditional probability, so the portfolio loses
\[
L(Z) = (1 - R)\, N\!\left(\frac{N^{-1}(p) - \sqrt{\rho}\,Z}{\sqrt{1 - \rho}}\right),
\]
and a \omterm{def:m2:credit-indices-and-tranches:tranche}{tranche}'s expected loss is the average of its loss over $Z$.

\begin{definition}[Base correlation]\label{def:m2:credit-indices-and-tranches:basecorr}
The \emph{base correlation}\index{base correlation} for a \omterm{def:m2:credit-indices-and-tranches:tranche}{detachment point} $d$ is the
correlation $\rho$ at which the one-factor model prices the \omterm{def:m2:credit-indices-and-tranches:tranche}{tranche} from 0 to $d$ at
its market value; a \omterm{def:m2:credit-indices-and-tranches:tranche}{tranche} from $a$ to $d$ is then valued as the difference of the
base \omterm{def:m2:credit-indices-and-tranches:tranche}{tranches} $[0,d]$ and $[0,a]$, each at its own correlation.
\end{definition}

Correlation moves risk between \omterm{def:m2:credit-indices-and-tranches:tranche}{tranches}
(\cref{fig:m2:credit-indices-and-tranches:tranches}). Low correlation makes defaults
independent: a few always happen, which hurts the equity, and many almost never, which
spares the senior. High correlation makes defaults come together or not at all, which
often spares the equity and sometimes hits the senior. So the equity is long
correlation and the senior short. The mezzanine sits between, and its expected loss
first rises and then falls with correlation, so one mezzanine price can fit two
correlations; \omterm{def:m2:credit-indices-and-tranches:basecorr}{base correlation}, defined on \omterm{def:m2:credit-indices-and-tranches:tranche}{tranches} that all start at zero, avoids
that ambiguity, since a base \omterm{def:m2:credit-indices-and-tranches:tranche}{tranche}'s expected loss always falls with correlation.

\begin{example}[Tranche losses]\label{ex:m2:credit-indices-and-tranches:tranches}
For the stylised index, the five-year default probability implied by the \omterm{def:m2:credit-indices-and-tranches:intrinsic}{intrinsic
spread} is 5.95\%, and the portfolio's expected loss 3.57\%. With $\rho = 0.3$ the
expected losses are 59.6\% of the equity \omterm{def:m2:credit-indices-and-tranches:tranche}{tranche}, 24.1\% of the mezzanine, 7.83\% of the
senior and 0.22\% of the super senior; weighted by their widths, 3, 4, 8 and 85\%, they
add up to 3.57\%. The equity's expected loss is 64.0\% at $\rho = 0.25$, and inverting
the model at that value returns a \omterm{def:m2:credit-indices-and-tranches:basecorr}{base correlation} of 25\%. The mezzanine's is 24.3\%
at $\rho = 0.10$, 25.0\% at 0.20 and 21.9\% at 0.45: it is 24.3\% again at 0.28.
\end{example}

\begin{omfigure}
\begin{tikzpicture}
\begin{axis}[width=11cm,height=5cm,
  xlabel={correlation $\rho$},ylabel={expected loss (\% of tranche)},
  tick label style={font=\small},label style={font=\small},
  xmin=0.05,xmax=0.9,ymin=0.001,ymax=100,ymode=log,grid=major,grid style={gray!25},
  legend columns=4,legend style={font=\footnotesize,at={(0.5,-0.3)},anchor=north,draw=none,fill=none,/tikz/every even column/.append style={column sep=8pt}},
  table/col sep=comma]
\addplot[omDef,very thick] table[x=rho,y=equity]{figdata/markets-2/24-credit-indices-and-tranches/tranches.csv};
\addplot[omThm,very thick] table[x=rho,y=mezz]{figdata/markets-2/24-credit-indices-and-tranches/tranches.csv};
\addplot[omProp,very thick] table[x=rho,y=senior]{figdata/markets-2/24-credit-indices-and-tranches/tranches.csv};
\addplot[gray,very thick,dashed] table[x=rho,y=super]{figdata/markets-2/24-credit-indices-and-tranches/tranches.csv};
\legend{0--3\%,3--7\%,7--15\%,15--100\%}
\end{axis}
\end{tikzpicture}

\omcaption{Five-year expected loss of each \omterm{def:m2:credit-indices-and-tranches:tranche}{tranche}, in per cent of its notional and on a
log scale, against the correlation of the one-factor Gaussian model (default
probability 5.95\%, recovery 40\%). The equity's expected loss falls with correlation
and the super senior's rises; the mezzanine's rises and then falls, and so, at high
correlations, does the 7--15\% \omterm{def:m2:credit-indices-and-tranches:tranche}{tranche}'s. Illustrative;
data: the chapter's tutorial.}\label{fig:m2:credit-indices-and-tranches:tranches}
\end{omfigure}

\begin{omfigure}
\begin{tikzpicture}
\begin{axis}[width=11cm,height=5cm,
  xlabel={portfolio loss (\% of notional)},ylabel={probability},
  tick label style={font=\small},label style={font=\small},
  xmin=-0.3,xmax=19.5,ymin=0,ymax=0.28,grid=major,grid style={gray!25},
  yticklabel style={/pgf/number format/.cd,fixed,precision=2},
  legend columns=2,legend style={font=\footnotesize,at={(0.5,-0.3)},anchor=north,draw=none,fill=none,/tikz/every even column/.append style={column sep=8pt}},
  table/col sep=comma]
\addplot[omDef,very thick,mark=*,mark size=1.1pt] table[x=loss,y=r10]{figdata/markets-2/24-credit-indices-and-tranches/losses.csv};
\addplot[omThm,very thick,mark=*,mark size=1.1pt] table[x=loss,y=r40]{figdata/markets-2/24-credit-indices-and-tranches/losses.csv};
\addplot[gray,thick,dashed,forget plot] coordinates {(3,0) (3,0.28)};
\addplot[gray,thick,dashed,forget plot] coordinates {(7,0) (7,0.28)};
\addplot[gray,thick,dashed,forget plot] coordinates {(15,0) (15,0.28)};
\legend{$\rho = 0.1$,$\rho = 0.4$}
\end{axis}
\end{tikzpicture}

\omcaption{Simulated distribution of the five-year loss of the 125-name portfolio, one
point per number of defaults (0.48\% each), for correlations of 0.1 and 0.4; the last
point collects 40 defaults or more. Dashed lines mark the 3, 7 and 15\% \omterm{def:m2:credit-indices-and-tranches:tranche}{attachment
points}. Higher correlation puts more weight on no defaults and on many. Illustrative;
data: the chapter's tutorial, 20\,000 seeded
draws.}\label{fig:m2:credit-indices-and-tranches:losses}
\end{omfigure}

The one-factor model is a convention for quoting, like the Black formula for
swaptions (chapter 13), not a description of how defaults happen; its failure to fit
all \omterm{def:m2:credit-indices-and-tranches:tranche}{tranches} with one correlation is the correlation skew: one number cannot describe
how defaults cluster.

\section{Index options}

\begin{definition}[Credit index option]\label{def:m2:credit-indices-and-tranches:option}
A \emph{credit index option}\index{credit index option} is an option to enter an
index contract at a fixed strike spread on an expiry date: a payer option gives the
right to buy protection, a receiver option the right to sell it.
\end{definition}

Index options trade on the five-year investment-grade and high-yield indices, European
style and settled by delivering the index position, with expiries of up to nine
months. A payer is worth exercising if the index spread at expiry is above the strike:
the holder buys protection at the strike's terms and receives the upfront difference,
roughly (spread minus strike) times the annuity. Payers are how investors buy
protection against a widening of \omterm{def:m2:corporate-bonds:spreads}{credit spreads} without paying the running premium,
and they are priced, in first approximation, with the Black formula on the spread,
the annuity playing the role that it plays for swaptions. The details, including how
defaults before expiry are treated and why the forward spread must allow for them, are
for One Quant Book 6.

\section{Tutorial: an index against its constituents, and its tranches}\label{tut:m2:credit-indices-and-tranches:tranche}

\begin{tutorial}
\textbf{Goal.} Compute a stylised index's \omterm{def:m2:credit-indices-and-tranches:intrinsic}{intrinsic spread} and skew from its 125
constituents, the P\&L of a skew trade, and \omterm{def:m2:credit-indices-and-tranches:tranche}{tranche} expected losses in the one-factor
Gaussian model, with a \omterm{def:m2:credit-indices-and-tranches:basecorr}{base correlation}.
\textbf{End state:} \cref{fig:m2:credit-indices-and-tranches:constituents,fig:m2:credit-indices-and-tranches:skewpnl,fig:m2:credit-indices-and-tranches:tranches,fig:m2:credit-indices-and-tranches:losses},
\cref{ex:m2:credit-indices-and-tranches:intrinsic,ex:m2:credit-indices-and-tranches:tranches}
and the numbers of the weekend problem.
\begin{tutsteps}
\item \textbf{The index}: intrinsic upfront, annuity-weighted spread, skew P\&L.
\omcode{code/firm/tranche/firm_tranche.py}{17}{31}{Intrinsic upfront, annuity-weighted spread and the P\&L of a skew position.}\label{lst:m2:credit-indices-and-tranches:index}
\item \textbf{The \omterm{def:m2:credit-indices-and-tranches:tranche}{tranches}}: \omterm{def:m2:credit-indices-and-tranches:tranche}{tranche} loss, large-pool expected loss, \omterm{def:m2:credit-indices-and-tranches:basecorr}{base correlation} and a
finite-pool simulation.
\omcode{code/firm/tranche/firm_tranche.py}{36}{73}{Tranche expected losses in the one-factor Gaussian model.}\label{lst:m2:credit-indices-and-tranches:tranche}
\item \textbf{Run} \texttt{index\_demo.tutorial()}, \texttt{index\_demo.skew\_trade()} and
\texttt{fig\_index.py}; the conversion between spreads and upfronts is chapter 23's
\texttt{firm.cds}.
\end{tutsteps}
\textbf{What to change next.} Make one constituent jump from 100 to 1\,000 basis points
and see how the \omterm{def:m2:credit-indices-and-tranches:intrinsic}{intrinsic spread} and the \omterm{def:m2:credit-indices-and-tranches:tranche}{tranches} move; then replace the large pool by
the finite 125-name simulation and measure the error in the equity \omterm{def:m2:credit-indices-and-tranches:tranche}{tranche}.
\end{tutorial}

\section{Build: index and tranche analytics}\label{bld:m2:credit-indices-and-tranches:tranche}

\begin{build}
\bfield{Purpose} The miniature firm marks its index positions against their
constituents, watches the skew, and prices \omterm{def:m2:credit-indices-and-tranches:tranche}{tranche} risk for its credit models.
\bfield{Interface} \texttt{intrinsic\_upfront(upfronts, weights)};
\texttt{annuity\_weighted\_spread(spreads, annuities)}; \texttt{skew\_pnl(notional,
upfront\_entry, upfront\_exit)}; \texttt{tranche\_loss(pool\_loss, attach, detach)};
\texttt{expected\_tranche\_loss(attach, detach, p, rho, recovery)};
\texttt{base\_correlation(detach, target, p, recovery)}; \texttt{simulate\_pool(names, p,
rho, recovery, trials, seed)}.
\bfield{Rules} Equal weights unless given; upfronts per unit notional from
\texttt{firm.cds}; large homogeneous pool with one default probability and one
recovery; expected losses by the trapezoid rule over the factor in $[-8, 8]$.
\bfield{Acceptance tests} \texttt{code/firm/tranche/tests/}: \omterm{def:m2:credit-indices-and-tranches:tranche}{tranche} losses clipped to
$[0, 1]$; \omterm{def:m2:credit-indices-and-tranches:tranche}{tranche} expected losses add up to the pool's; equity falls and senior rises
with correlation; \omterm{def:m2:credit-indices-and-tranches:basecorr}{base correlation} round-trips; a finite 125-name pool agrees with the
large pool.
\bfield{Stretch} Heterogeneous default probabilities with the recursion for the
finite pool's loss distribution; stochastic recovery; bootstrapping \omterm{def:m2:credit-indices-and-tranches:basecorr}{base correlations}
from \omterm{def:m2:credit-indices-and-tranches:tranche}{tranche} quotes; index option pricing (One Quant Book 6).
\end{build}

\begin{omsources}
\item SEC, Order Instituting Cease-and-Desist Proceedings, In the Matter of JPMorgan
Chase \& Co., Release No.~34-70458, 19 September 2013.
\item JPMorgan Chase \& Co., Report of the Management Task Force Regarding 2012 CIO
Losses, 16 January 2013.
\item US Senate, Permanent Subcommittee on Investigations, hearing ``JPMorgan Chase
whale trades: a case history of derivatives risks and abuses'', 15 March 2013.
\item ICE, Markit CDX.NA.IG contract specifications; CDS index options clearing.
\item Bloomberg SEF, CDX.NA.IG five-year contract submission to the CFTC, November 2013.
\end{omsources}

\section{Exercises}

\begin{exercise}[$\star$]\label{exo:m2:credit-indices-and-tranches:1}
You sold USD~1 billion of protection on a 125-name index. One name defaults and its
\omterm{def:m2:credit-default-swaps:final}{auction final price} is 40. What do you pay, and what is the index notional afterwards?
\end{exercise}

\begin{exercise}[$\star$]\label{exo:m2:credit-indices-and-tranches:2}
What fraction of the 3--7\% \omterm{def:m2:credit-indices-and-tranches:tranche}{tranche} is lost if portfolio losses are 2\%, 5\% and 9\%?
\end{exercise}

\begin{exercise}[$\star$]\label{exo:m2:credit-indices-and-tranches:3}
Why does a new \omterm{def:m2:credit-indices-and-tranches:index}{index series} draw most of the trading, and what happens to the old one?
\end{exercise}

\begin{exercise}[$\star\star$]\label{exo:m2:credit-indices-and-tranches:4}
An index has two names, at 50 and 250 basis points. Is its \omterm{def:m2:credit-indices-and-tranches:intrinsic}{intrinsic spread} 150? Use
the chapter's conventions.
\end{exercise}

\begin{exercise}[$\star\star$]\label{exo:m2:credit-indices-and-tranches:5}
Why is a seller of \omterm{def:m2:credit-indices-and-tranches:tranche}{equity-tranche} protection long correlation?
\end{exercise}

\begin{exercise}[$\star\star$]\label{exo:m2:credit-indices-and-tranches:6}
An investor buys a three-month payer option on the index struck at 80 basis points.
When does it exercise, and what does it get?
\end{exercise}

\begin{exercise}[$\star\star\star$]\label{exo:m2:credit-indices-and-tranches:7}
\emph{Coding.} With \texttt{expected\_tranche\_loss}, check that the four \omterm{def:m2:credit-indices-and-tranches:tranche}{tranches}'
expected losses add up to the portfolio's at $\rho = 0.3$, and find the second
correlation at which the mezzanine's expected loss equals its value at $\rho = 0.1$.
\end{exercise}

\begin{exercise}[$\star\star\star$]\label{exo:m2:credit-indices-and-tranches:8}
\emph{Find the flaw.} ``The index is the average of its 125 spreads. It trades 8 basis
points below that average, so buy index protection and sell the names: a riskless
profit.'' Correct it.
\end{exercise}

\section{Problem: The Skew Trade}

\begin{problem}[{Weekend problem --- trading an index against its constituents}]\label{pb:m2:credit-indices-and-tranches:1}
A relative-value fund watches the stylised index of
\cref{ex:m2:credit-indices-and-tranches:intrinsic}: 125 names, \omterm{def:m2:credit-indices-and-tranches:intrinsic}{intrinsic spread} 73.6
basis points, while a large seller of index protection has pushed the index to 65.6.
The fund buys USD~1 billion of index protection and sells USD~8 million of protection
on each constituent. Trading a single name costs half a bid--offer spread of 1.5 basis
points of running spread, the index 0.25; flat rate 4\%, recovery 40\%.

\medskip\noindent\textbf{Part I --- The index and its constituents.}
\begin{enumerate}
\item What are the simple average and the median of the constituents' spreads?
\item What is the intrinsic upfront, and the \omterm{def:m2:credit-indices-and-tranches:intrinsic}{intrinsic spread}?
\item Why is the \omterm{def:m2:credit-indices-and-tranches:intrinsic}{intrinsic spread} below the simple average?
\item Compare it with the annuity-weighted average of the spreads.
\item What is the skew?
\end{enumerate}

\medskip\noindent\textbf{Part II --- The trade.}
\begin{enumerate}[resume]
\item What upfront does the fund receive on the index, and pay on the constituents?
\item What is the net cash at entry, and what does it represent?
\item One constituent defaults with an \omterm{def:m2:credit-default-swaps:final}{auction final price} of 40. What does the fund pay
and receive?
\item What is the P\&L if the skew closes, constituents unchanged?
\item What is the P\&L if the skew first widens to $-20$ basis points?
\end{enumerate}

\medskip\noindent\textbf{Part III --- Costs and risks.}
\begin{enumerate}[resume]
\item What does it cost to enter both legs, and to enter and leave them?
\item What is the P\&L net of costs if the skew closes, and by how much must the skew
move for the trade to pay its costs?
\item Why can the skew persist or widen?
\item What does a widening cost the fund before it ends, beyond the loss on paper?
\item Why is a very large seller of index protection exposed to skew traders?
\end{enumerate}

\medskip\noindent\textbf{Part IV --- Judgement.}
\begin{enumerate}[resume]
\item Why could the seller not simply unwind a position of 10 to 15 days of the
market's volume?
\item How would you size a skew trade?
\item What happens to the trade at the next roll?
\item State the \emph{named result}: the P\&L of the skew trade when the skew closes,
gross and net of costs, and its loss if the skew widens first.
\item In one sentence: why is the skew not a free lunch?
\end{enumerate}
\end{problem}

\section{Interview questions}

\begin{interviewq}[$\star$ \iqroles{trader, researcher}]\label{iq:m2:credit-indices-and-tranches:1}
What is CDX investment grade, and how does it roll?
\end{interviewq}

\begin{interviewq}[$\star$ \iqroles{trader}]\label{iq:m2:credit-indices-and-tranches:2}
What is an index's \omterm{def:m2:credit-indices-and-tranches:intrinsic}{intrinsic spread}, and why can the index trade away from it?
\end{interviewq}

\begin{interviewq}[$\star\star$ \iqroles{researcher}]\label{iq:m2:credit-indices-and-tranches:3}
Derive the loss of a large pool in the one-factor Gaussian model.
\end{interviewq}

\begin{interviewq}[$\star\star$ \iqroles{researcher, trader}]\label{iq:m2:credit-indices-and-tranches:4}
Which \omterm{def:m2:credit-indices-and-tranches:tranche}{tranches} are long correlation, and why? Why use \omterm{def:m2:credit-indices-and-tranches:basecorr}{base correlation} rather than
one correlation per \omterm{def:m2:credit-indices-and-tranches:tranche}{tranche}?
\end{interviewq}

\begin{interviewq}[$\star\star$ \iqroles{risk, trader}]\label{iq:m2:credit-indices-and-tranches:5}
What went wrong in the 2012 losses of JPMorgan's Chief Investment Office?
\end{interviewq}

\begin{interviewq}[$\star\star\star$ \iqroles{developer}]\label{iq:m2:credit-indices-and-tranches:6}
Design a service that computes, in real time, the \omterm{def:m2:credit-indices-and-tranches:intrinsic}{intrinsic spread} and skew of every
\omterm{def:m2:credit-indices-and-tranches:index}{index series} from single-name quotes.
\end{interviewq}
